\section[Notebook 1968-1: a simple module is determined by its trace]{Notebook 1968-1: a simple module is determined by its trace\protect\verifie{Notebook1968_1}}\label{sec:quillen-1968-1}

\emph{Context.} \carnet{1968}{1}, dated January 1968, is mostly about representations of semisimple Lie algebras. On PDF page~4, Quillen isolates a basic statement: an irreducible representation of a finite-dimensional algebra over $\mathbb{C}$ is determined by its trace function $a \mapsto \operatorname{tr}(a \mid S)$. His sketch: the trace vanishes on the radical, whose elements act nilpotently; by Wedderburn, the semisimple quotient is a product of simple algebras, and the trace of $S$ is nonzero on the factor of $S$ and zero on the others. A remark that the trace equals $\dim S$ at $1$ is crossed out, rightly so in characteristic $p$, where $\dim S$ can vanish in the field.

\emph{Intuition.} If $S \not\simeq T$, it suffices to find an element of $A$ acting on $S$ as a projector of rank one and on $T$ as zero: its traces are $1$ and $0$, in every characteristic. Jacobson's density theorem provides such an element, because the endomorphisms of $S \times T$ are diagonal and scalar.

\begin{encadre}
\begin{theoreme}[Quillen, notebook 1968-1, p.~4]\label{thm:quillen-1968-1}
Let $k$ be an algebraically closed field of arbitrary characteristic, $A$ a $k$-algebra, and $S$, $T$ two simple $A$-modules of finite dimension over $k$. If $\operatorname{tr}(a \mid S) = \operatorname{tr}(a \mid T)$ for all $a \in A$, then $S \simeq T$. More precisely, if $S \not\simeq T$, there exists $a \in A$ with $\operatorname{tr}(a \mid S) = 1$ and $\operatorname{tr}(a \mid T) = 0$. When $A$ is finite-dimensional, as in the notebook, the finite-dimensionality hypothesis on $S$ and $T$ is automatic.
\end{theoreme}
\end{encadre}

\begin{lemme}[Schur]\label{lem:quillen-1968-1-schur}
An $A$-linear map between non-isomorphic simple modules is zero. If $k$ is algebraically closed and $S$ is of finite dimension over $k$, every $A$-linear endomorphism of $S$ is a scalar multiplication.
\end{lemme}

\begin{proof}
The kernel and the image are submodules, so a nonzero map between simple modules is bijective. An endomorphism $g$ of $S$ has an eigenvalue $c \in k$; $g - c$ is $A$-linear and not injective, hence zero.
\end{proof}

\begin{proof}[Proof of the theorem]
The module $S \times T$ is semisimple, being the sum of the images of the two injections. By Lemma~\ref{lem:quillen-1968-1-schur}, if $S \not\simeq T$, an $A$-linear endomorphism $g$ of $S \times T$ has the form $g(s, t) = (c\,s, d\,t)$ with $c, d \in k$. Let $v \neq 0$ in $S$, $\varphi$ a linear form with $\varphi(v) = 1$, and $\psi = \varphi(\cdot)\,v$, of trace $1$. The map $f(s, t) = (\psi(s), 0)$ commutes with all these $g$; since $S \times T$ is finitely generated over its endomorphism ring (it already is over $k$), the density theorem gives $a \in A$ acting as $f$. Then $\operatorname{tr}(a \mid S) = \operatorname{tr}\psi = 1$ and $\operatorname{tr}(a \mid T) = 0$, which contradicts the equality of traces. If $A$ is finite-dimensional, a simple module is cyclic, hence finite-dimensional over $k$.
\end{proof}

\begin{remarque}
The proof avoids the radical and Wedderburn--Artin, which Quillen uses; the two routes give the same element, a matrix unit $e_{11}$ of the block of $S$. The finiteness of $A$ is not used, only that of $S$ and $T$.
\end{remarque}

\begin{contrexemple}
The hypothesis \q{algebraically closed} cannot be dropped in characteristic $p$. Let $L/k$ be a purely inseparable extension of degree $p$ and $A = L \times L$; the two simple modules $L \times 0$ and $0 \times L$ are not isomorphic, and their traces, $(x, y) \mapsto \operatorname{Tr}_{L/k}(x)$ and $(x, y) \mapsto \operatorname{Tr}_{L/k}(y)$, are both zero. Likewise, over an algebraically closed field of characteristic $p$, the trace does not determine semisimple modules: $S^{\oplus p}$ and the zero module have the same trace.
\end{contrexemple}

\begin{sloppypar}
\noindent\emph{Verification in Lean.} The theorem is \texttt{nonempty\_linearEquiv\_of\_trace\_eq}, and \texttt{nonempty\_linearEquiv\_of\_trace\_eq'} for $A$ finite-dimensional; the element with traces $1$ and $0$ is \texttt{exists\_trace\_eq\_one\_trace\_eq\_zero}. Lemma~\ref{lem:quillen-1968-1-schur} is \texttt{eq\_zero\_of\_isEmpty\_linearEquiv} and \texttt{exists\_eq\_smul\_of\_isAlgClosed}, the semisimplicity of $S \times T$ is \texttt{isSemisimpleModule\_prod}, in \href{\leanblobmain/Notebook1968_1.lean}{\texttt{Notebook1968\_1.\allowbreak{}lean}}. Density is mathlib's \texttt{Module.Finite.toModuleEnd\_moduleEnd\_surjective}.
\end{sloppypar}

Not formalized: the counterexample above, and the linear independence of the traces of pairwise non-isomorphic simple modules, which follows from the same argument.

\litterature{found in the literature: Curtis--Reiner, \emph{Representation Theory of Finite Groups and Associative Algebras} (1962), the chapter on characters of algebras; Lam, \emph{A First Course in Noncommutative Rings}, the chapter on representations. mathlib has the density theorem and Wedderburn--Artin over an algebraically closed field, but not this statement; it could be proposed for mathlib, and the edition does not present it as new.}
